% Copyright (c) 2026 Eugene Joseph M. Ragasa.
% SPDX-License-Identifier: Apache-2.0
% Author: Eugene Joseph M. Ragasa
% Formal proof completed: 2026-10-20
\documentclass[11pt]{article}
\usepackage[T1]{fontenc}\usepackage{amsmath,amssymb}\usepackage[margin=1in]{geometry}\usepackage[hidelinks]{hyperref}
\title{Proof Dossier: Finite First Joint Pass}\author{Eugene Joseph M. Ragasa}\date{Formal proof completed: 2026-10-20}
\begin{document}\maketitle
\section{Identity and status}
ID \texttt{BG-DECISION-THM-016}; status \texttt{proved}. Lean declaration:
\par\noindent{\tiny\path{ProjectKoios.BandGauge.Decision.}}\\
{\tiny\path{isFirstJointPassing_iff}}.
\section{Statement and proof}
A first joint pass is exactly one inventory cutoff that passes both predicates while
every strictly earlier inventory cutoff fails at least one. The equivalence unfolds the
first-pass definition and reassociates conjunctions.
\section{Boundary}
Separate channel witnesses do not establish this proposition. The theorem does not bind
a runtime table or prove inventory completeness. Kernel axioms are \texttt{propext},
\texttt{Classical.choice}, and \texttt{Quot.sound}; no admissions. Checked with Lean
\texttt{4.35.0} and Mathlib revision
\texttt{56a005855ea17e0b0958f73ec801ae140b0a1679}. AI-assisted and kernel-checked;
no independent expert review recorded.
\end{document}
